\documentclass[../main.tex]{subfiles}

\begin{document}

\ifSubfilesClassLoaded{

    \tableofcontents
    
}{}

\chapter{ Continuous Functions }

\section{Continuity of a Function}


\begin{definition}{}{}
Let $X$ and $Y$ be topological spaces. A function $f: X \to Y$ is said to be \dfntxt{continuous} if for each open subset $V$ of $Y$, the set $f^{-1}(V)$ is an open subset of $X$.
\end{definition}

\begin{theorem}{}{}
Let $X$ and $Y$ be topological spaces; let $f : X \to Y$. Then the following are equivalent:
\begin{enumerate}
    \item $f$ is continuous.
    \item For every subset $A$ of $X$, one has $f(\bar{A}) \subset \overline{f(A)}$.
    \item For every closed set $B$ of $Y$, the set $f^{-1}(B)$ is closed in $X$.
    \item For each $x \in X$ and each neighborhood $V$ of $f(x)$, there is a neighborhood $U$ of $x$ such that $f(U) \subset V$.
\end{enumerate}
\end{theorem}
\begin{remark}
    
If the condition in (4) holds for the point $x$ of $X$, we say that $f$ is \dfntxt{continuous at the point x}.

\end{remark}

\begin{proof}
    1. \(  \implies   \) 2. : If \(  f  \) is continuous, take any  \(  x \in f\left( \bar{A} \right)   \). Let \(  V  \) be any neighborhood of \(  f\left( x \right)   \). Then \(  f^{-1} \left( V \right)   \) is a neighborhood of \(  x  \). We have \(  f^{-1} \left( V \right)\cap A\neq \varnothing   \). Then \[
    \varnothing\neq  f\left( f^{-1} \left( V \right)\cap A  \right)\subseteq V\cap f\left( A \right)  
    \]        Since \(  V  \) is arbitrary, \(  f\left( x \right) \in  \overline{f\left( A \right) }   \).  
\end{proof}

\section{Homeomorphisms}

\begin{definition}{}{}
    Let $X$ and $Y$ be topological spaces; let $f: X \to Y$ be a bijection. If both the function $f$ and the inverse function
\[ f^{-1} : Y \to X \]
are continuous, then $f$ is called a \dfntxt{homeomorphism}.
\end{definition}

\end{document}